\documentclass[9.5pt,a4paper]{article}

\usepackage[a4paper,top=12mm,bottom=12mm,left=13mm,right=13mm]{geometry}
\usepackage{fontspec}
\usepackage{xcolor}
\usepackage{tabularx}
\usepackage{array}
\usepackage[most]{tcolorbox}
\usepackage{enumitem}
\usepackage{ragged2e}

\IfFontExistsTF{Arial}{
  \setmainfont{Arial}
  \setsansfont{Arial}
}{
  \setmainfont{TeX Gyre Heros}
  \setsansfont{TeX Gyre Heros}
}

\definecolor{AIEPRed}{HTML}{D62027}
\definecolor{AIEPNavy}{HTML}{082B5C}
\definecolor{AIEPInk}{HTML}{243B53}
\definecolor{AIEPMuted}{HTML}{5F6B7A}
\definecolor{AIEPLine}{HTML}{CBD5E1}
\definecolor{AIEPBlueSoft}{HTML}{E9F1F8}
\definecolor{AIEPGreen}{HTML}{2E8B57}
\definecolor{AIEPGold}{HTML}{D8A928}

\pagestyle{empty}
\setlength{\parindent}{0pt}
\setlength{\parskip}{2.5pt}
\hyphenpenalty=10000
\exhyphenpenalty=10000
\setlist[itemize]{leftmargin=4.4mm,topsep=1pt,itemsep=.8pt,parsep=0pt}
\renewcommand{\arraystretch}{1.14}
\newcolumntype{Y}{>{\RaggedRight\arraybackslash}X}
\newcolumntype{L}[1]{>{\RaggedRight\arraybackslash}p{#1}}

\newtcolorbox{boxa}[2][]{
  enhanced,colback=white,colframe=#2,boxrule=.55pt,arc=1mm,
  left=3mm,right=3mm,top=2.2mm,bottom=2.2mm,#1
}
\newcommand{\labelx}[1]{{\scriptsize\bfseries\color{AIEPRed}\MakeUppercase{#1}}\par}

\begin{document}

\begin{tcolorbox}[colback=AIEPNavy,colframe=AIEPNavy,arc=1.5mm,left=5mm,right=5mm,top=4mm,bottom=4mm]
{\small\bfseries\color{white}GEOGREEN ESCOLAR · AIEP OSORNO}\par
{\LARGE\bfseries\color{white}Taller 3 · De la idea al prototipo}\par
{\normalsize\color{white}Guía rápida para ejecutar el paquete docente}
\end{tcolorbox}

\begin{tabularx}{\textwidth}{Y Y Y}
\begin{boxa}[height=28mm,valign=top]{AIEPRed}
\labelx{Cuándo}
\textbf{Lunes 24 de agosto}\par
Bloque A: 09:00–10:30\par
Bloque B: 10:45–12:15
\end{boxa}
&
\begin{boxa}[height=28mm,valign=top]{AIEPGold}
\labelx{Organización}
\textbf{30 estudiantes por bloque}\par
5 equipos de 6 integrantes
\end{boxa}
&
\begin{boxa}[height=28mm,valign=top]{AIEPGreen}
\labelx{Responsabilidad}
\textbf{Programación y Análisis de Sistemas}\par
Sesión de 90 minutos
\end{boxa}
\end{tabularx}

\begin{boxa}{AIEPRed}
\labelx{Propósito}
Que cada equipo comprenda cómo una necesidad puede convertirse en un sistema tecnológico y salga con
una propuesta inicial propia: problema, variable, sensor, regla, prueba segura y siguiente hito.
GeoGreen funciona como caso inspirador y demuestra el ciclo \textbf{sensar → procesar → comunicar → actuar}.
\end{boxa}

\begin{tabularx}{\textwidth}{Y Y}
\begin{boxa}[height=61mm,valign=top]{AIEPNavy}
\labelx{Mapa de la clase}
\begin{itemize}
  \item \textbf{0–15 min · Una idea puede crecer.} GeoGreen, trayectoria y desafío.
  \item \textbf{15–35 min · Del mundo físico al dato.} Placa, protoboard, sensores, actuadores y seguridad.
  \item \textbf{35–70 min · Prototipar con sensores e IA.} Variable, sensor, regla, simulación o primera prueba.
  \item \textbf{70–90 min · Del prototipo al producto.} Madurez, software, carcasa, PCB y siguiente hito.
\end{itemize}
\end{boxa}
&
\begin{boxa}[height=61mm,valign=top]{AIEPGreen}
\labelx{Salida obligatoria}
Cada equipo entrega la \textbf{Ficha de propuesta tecnológica inicial} con:
\begin{itemize}
  \item problema, contexto y personas afectadas;
  \item variable y sensor justificado;
  \item regla “cuando..., entonces...”;
  \item evidencia inicial o prueba pendiente;
  \item revisión de seguridad;
  \item seis responsabilidades y próximo hito.
\end{itemize}
\end{boxa}
\end{tabularx}

\begin{boxa}{AIEPGold}
\labelx{Material incluido}
\begin{tabularx}{\textwidth}{@{}L{43mm}Y@{}}
\textbf{Planificación docente} & Secuencia completa, mediación, seguridad, evaluación y continuidad.\\
\textbf{Presentación} & Apoyo visual de 67 diapositivas con videos integrados.\\
\textbf{Ficha de equipo} & Documento imprimible; cinco copias por bloque.\\
\textbf{Mapa e infografías} & Resumen de la clase y una lámina vertical por bloque.\\
\textbf{Contexto GeoGreen} & Resumen del proyecto, lógica general y roles de los equipos.\\
\end{tabularx}
\end{boxa}

\begin{tabularx}{\textwidth}{Y Y}
\begin{boxa}[height=42mm,valign=top]{AIEPRed}
\labelx{Antes de comenzar}
\begin{itemize}
  \item Probar el PPT y sus dos videos.
  \item Disponer cinco estaciones de exploración.
  \item Revisar sensores, cables y alimentación.
  \item Imprimir cinco fichas por bloque.
  \item Dejar el prototipo GeoGreen listo para demostrar.
\end{itemize}
\end{boxa}
&
\begin{boxa}[height=42mm,valign=top]{AIEPNavy}
\labelx{Regla de seguridad}
La placa permanece desconectada durante el armado. Antes del USB se revisan pinout, voltaje,
polaridad, resistencias y tierra común. La IA puede proponer; el equipo y quien facilita deben
verificar antes de conectar o ejecutar.
\end{boxa}
\end{tabularx}

\begin{tcolorbox}[colback=AIEPBlueSoft,colframe=AIEPNavy,boxrule=.65pt,arc=1mm,left=3mm,right=3mm,top=2.5mm,bottom=2.5mm]
\labelx{Criterio de ejecución}
La demostración abre posibilidades; no reemplaza el trabajo del equipo. Proteger el tramo práctico,
recoger todas las fichas y dejar claro el avance esperado antes de la Mentoría 1.
\end{tcolorbox}

\begin{tcolorbox}[colback=AIEPRed,colframe=AIEPRed,arc=1mm,halign=center,top=2.5mm,bottom=2.5mm]
{\large\bfseries\color{white}GeoGreen demuestra hasta dónde puede crecer una idea. Hoy comienza la de ustedes.}
\end{tcolorbox}

\end{document}
